
\subsection{Finite Trotter slices and a two-site transfer matrix}
\label{sm:sec:transfer}

OpenAI's proof uses a continuous-time spatial integral kernel~\cite{OpenAI2026Periodic}.
For our Lean formalization we instead prove the spatial inequalities
for finite matrices and pass them to exact partition functions by a
controlled product limit.
% This section gives that alternative proof.
% The subsequent doubling argument is the original one.

Consider the periodic spin-1 Heisenberg Hamiltonian
\begin{equation}
 H_L=\sum_{j=0}^{L-1}h_{j,j+1},
 \qquad h=\sum_{\alpha=x,y,z}S^\alpha\otimes S^\alpha,
 \qquad Z_L(\beta)=\Tr e^{-\beta H_L},
 \label{sm:eq:model}
\end{equation}
with indices taken modulo $L$, exchange coupling equal to one, and
$\beta>0$. In the Cartesian basis $s,s'\in\{x,y,z\}$, set
\begin{equation}
 (S^\alpha)_{ss'}=-i\epsilon_{\alpha ss'},
 \qquad G_\alpha=iS^\alpha,
 \qquad -h=\sum_\alpha G_\alpha\otimes G_\alpha.
 \label{sm:eq:generators}
\end{equation}
Here $\epsilon_{\alpha ss'}$ is the Levi-Civita symbol, $\Tr$ denotes
the matrix trace, and $\smI$ denotes the identity on the relevant space.
Each $G_\alpha$ is real antisymmetric and has operator norm one.
Equation~\eqref{sm:eq:generators} is the real bond-reflection-positive
factorization used below~\cite{Dyson1978Phase,Frohlich1978Phase}.

To replace the continuous-time path integral by a finite-dimensional
trace, we discretize imaginary time into $M$ brickwork slices.
For $L=2N$, $N\ge1$, and a positive integer $M$, let
$\tau=\beta/M$. Define the linearized brickwork slice and its trace by
\begin{align}
 S_{2N}(\tau)
 &=\prod_{k=0}^{N-1}(\smI-\tau h_{2k,2k+1})
   \prod_{k=0}^{N-1}(\smI-\tau h_{2k+1,2k+2}), \notag\\
 \smZ{2N}(\beta)&=\Tr S_{2N}(\beta/M)^M.
 \label{sm:eq:slice}
\end{align}
Products are ordered from left to right in increasing index.
The trace is real because all factors have real entries.
% We do not assume that the individual linearized factors are positive.
% mannered prose here.

Let $\Gamma=\{\varnothing,x,y,z\}$, with
$C_\varnothing=\smI$ and $C_\alpha=\sqrt{\tau}\,G_\alpha$.
The bond expansion is the exact finite identity
\begin{equation}
 \smI-\tau h=\sum_{\gamma\in\Gamma}C_\gamma\otimes C_\gamma.
 \label{sm:eq:bond}
\end{equation}
An imaginary-time bond sequence is an element
$p=(p_0,\ldots,p_{M-1})$ of $\smD=\Gamma^M$.
Define the single-site trace matrix by
\begin{equation}
 \smF_{pq}
 =\Tr\!\left[\prod_{t=0}^{M-1}(C_{p_t}C_{q_t})\right],
 \qquad p,q\in\smD.
 \label{sm:eq:F}
\end{equation}
The product over $t$ follows the multiplication order of the $M$ slices.
The matrix $\smF$ depends on $M$ and $\tau$, but not on $N$.

To contract the network in space, assign sequences $b_k$ and $c_k$
to the even and odd bonds starting at $2k$ and $2k+1$, respectively.
In slice $t$, site $2k$ carries
$C_{b_{k,t}}C_{c_{k-1,t}}$, whereas site $2k+1$ carries
$C_{b_{k,t}}C_{c_{k,t}}$. Expanding all bond sums, multiplying
sitewise tensor products, and taking the physical trace gives
\begin{align}
 \smZ{2N}(\beta)
 &=\sum_{\substack{b_0,\ldots,b_{N-1}\in\smD\\
                   c_0,\ldots,c_{N-1}\in\smD}}
   \prod_{k=0}^{N-1}
   \smF_{b_kc_{k-1}}\smF_{b_kc_k} \notag\\
 &=\sum_{c_0,\ldots,c_{N-1}\in\smD}
   \prod_{k=0}^{N-1}(\smF^{\mathsf T}\smF)_{c_{k-1}c_k}
   =\Tr(\smF^{\mathsf T}\smF)^N,
 \label{sm:eq:factorization}
\end{align}
where $c_{-1}=c_{N-1}$. The middle equality sums each even-bond
sequence $b_k$ independently; the last equality is the closed-path
expansion of a matrix trace.

Thus the spatial transfer matrix
\begin{equation}
 \smT=\smF^{\mathsf T}\smF=\smF^\dagger\smF\succeq0
 \label{sm:eq:positive}
\end{equation}
advances \emph{two sites}, not one. Because $\smF$ is real, its
transpose equals its adjoint, and
$v^\dagger\smT v=\smnorm{\smF v}^2\ge0$.
Neither symmetry nor positivity of $\smF$ itself is required.
This finite two-site Gram matrix is distinct from the
continuous-time one-site self-adjoint kernel used in a thermal-MPO
description~\cite{Tang2021Continuous}, and we do not construct a
limiting kernel.

We extend the same contraction to any real reflection-positive cell model.
Its assumptions are a finite-dimensional local cell space, a finite
family of real matrices $C_\alpha$, each symmetric or antisymmetric,
a real symmetric on-site matrix $o$, and a real symmetric involution
$F$ satisfying
\begin{equation}
 F^2=\smI,\qquad FoF=o,\qquad
 h=-\sum_\alpha C_\alpha\otimes FC_\alpha F.
 \label{sm:eq:cell}
\end{equation}
The Hamiltonian is $H_L=\sum_j h_{j,j+1}+\sum_j o_j$, and lengths
count cells. Insert the on-site layer
$\prod_i(\smI-\tau o_i)$ between the even and odd bond layers.
Writing $A_\varnothing=\smI$,
$A_\alpha=\sqrt{\tau}\,C_\alpha$, and $O=\smI-\tau o$, the bond
expansion becomes
$\smI-\tau h=\sum_\gamma A_\gamma\otimes FA_\gamma F$.
Replace Eq.~\eqref{sm:eq:F} by
\begin{equation}
 \smF_{pq}=\Tr\prod_{t=0}^{M-1}(A_{p_t}OFA_{q_t}F),
 \label{sm:eq:cell-F}
\end{equation}
where the sequences now range over the enlarged bond-label alphabet.
An odd cell carries $FA_{p_t}FOA_{q_t}$. Since $FO=OF$,
this factor is the conjugate by $F$ of the even-cell factor.
Conjugating the entire product leaves its trace unchanged.
The two sublattices therefore contribute the same $\smF_{pq}$,
and Eqs.~\eqref{sm:eq:factorization} and~\eqref{sm:eq:positive}
follow without change.

\subsection{Finite-transfer moment inequalities}
\label{sm:sec:moments}

Positivity of $\smT$ expresses $\smZ{L}(\beta)$ as an even moment sum of
nonnegative weights, from which we derive the spatial purity, squaring,
and interpolation inequalities at finite $M$.
At fixed $M,\beta$, let $\lambda_i\ge0$ be the eigenvalues of
$T_{M,\beta/M}$, and put $\mu_i=\sqrt{\lambda_i}$.
For every even $L\ge2$, Eq.~\eqref{sm:eq:factorization} gives
\begin{equation}
 \smZ{L}(\beta)=\sum_i\mu_i^L.
 \label{sm:eq:moments}
\end{equation}

For any finite nonnegative family $a_i$, let $S_k=\sum_i a_i^k$.
Then
\begin{equation}
 S_2\le S_1^2,\qquad
 2(S_2^2-S_4)\le(S_1^2-S_2)^2.
 \label{sm:eq:powers}
\end{equation}
Indeed, $S_1^2-S_2=2\sum_{i<j}a_i a_j$, whereas
$S_2^2-S_4=2\sum_{i<j}a_i^2a_j^2$.
Squaring the first sum and retaining only its diagonal terms proves
the second inequality.

If $\sum_i a_i=1$, the purity is $P=\sum_i a_i^2>0$.
Squaring and normalizing replaces the distribution by
$a_i^2/P$, of purity $P'=S_4/P^2$. Thus
\begin{equation}
 1-P'\le f(1-P),\qquad
 f(z)=\frac{z^2}{2(1-z)^2}.
 \label{sm:eq:squaring}
\end{equation}
The map $f$ is nondecreasing on $0\le z<1$, since
$z/(1-z)$ is nonnegative and increasing there.

For $0<u\le v$, nonnegative weights also satisfy
\begin{equation}
 \sum_i\mu_i^v\le\left(\sum_i\mu_i^u\right)^{v/u}.
 \label{sm:eq:lp}
\end{equation}
If the sum on the right is zero, the claim is immediate.
Otherwise normalize $\mu_i^u$ to probabilities $a_i$; then
$\sum_i a_i^{v/u}\le\sum_i a_i=1$.

Applying Eqs.~\eqref{sm:eq:powers} and~\eqref{sm:eq:lp} to
Eq.~\eqref{sm:eq:moments} yields, at fixed $\beta$, the closed scalar
inequalities
\begin{align}
 \smZ{2n}&\le(\smZ{n})^2, \notag\\
 2((\smZ{2n})^2-\smZ{4n})
 &\le((\smZ{n})^2-\smZ{2n})^2, \notag\\
 \smZ{v}&\le(\smZ{u})^{v/u},
 \label{sm:eq:closed}
\end{align}
for even $n\ge2$ and even $v\ge u\ge2$.
We suppress the temperature argument only in this display.

\subsection{Controlled product limit}
\label{sm:sec:limit}

To pass the finite-$M$ inequalities in Eq.~\eqref{sm:eq:closed} to the
exact partition functions $Z_L(\beta)$, we prove the matrix product
limit $S_L(\beta/M)^M\to e^{-\beta H_L}$ at fixed $L$ and $\beta$.
Use the operator norm, which is submultiplicative and satisfies
$\smnorm{\smI}=1$.
For $\smnorm{X},\smnorm{Y}\le K$, $K\ge1$, telescoping powers
gives
\begin{equation}
 \smnorm{X^k-Y^k}\le kK^k\smnorm{X-Y}.
 \label{sm:eq:telescoping}
\end{equation}
For $k\ge1$, expand the difference as
$\sum_{r=0}^{k-1}X^{k-1-r}(X-Y)Y^r$; each term is bounded by
$K^{k-1}\smnorm{X-Y}\le K^k\smnorm{X-Y}$.
The case $k=0$ is immediate.

Suppose $M(X_M-\smI)\to A$, and set $Y_M=e^{A/M}$.
Then $M(Y_M-\smI)\to A$ and $Y_M^M=e^A$.
With $C=\smnorm{A}+1$, convergence implies, for sufficiently large
$M$,
\[
 \smnorm{X_M},\smnorm{Y_M}\le1+C/M.
\]
Equation~\eqref{sm:eq:telescoping} now gives the quantitative comparison
\begin{align}
 \smnorm{X_M^M-e^A}
 &\le (1+C/M)^M\smnorm{M(X_M-Y_M)} \notag\\
 &\le e^C\smnorm{M(X_M-Y_M)}\longrightarrow0.
 \label{sm:eq:error}
\end{align}
In particular, if $g(0)=\smI$ and $g'(0)=A$, then
$g(1/M)^M\to e^A$.

For fixed $L,\beta$, list the negated bond operators in the order
of Eq.~\eqref{sm:eq:slice}, including the negated on-site operators
when present, as $B_1,\ldots,B_k$.
Their sum is $-H_L$. The polynomial
\[
 g(t)=\prod_{i=1}^k(\smI+\beta tB_i)
\]
has derivative $g'(0)=-\beta H_L$. Hence
\begin{equation}
 S_L(\beta/M)^M\longrightarrow e^{-\beta H_L},
 \qquad \smZ{L}(\beta)\longrightarrow Z_L(\beta).
 \label{sm:eq:limit}
\end{equation}
More explicitly, Taylor expansion gives
$g(1/M)-e^{-\beta H_L/M}=O(M^{-2})$.
Equation~\eqref{sm:eq:error} therefore bounds the operator error by
$e^C K_{L,\beta}/M$ for some finite $K_{L,\beta}$ and sufficiently
large $M$, with $C=\smnorm{\beta H_L}+1$.
Trace convergence follows from
$\lvert\Tr A\rvert\le3^L\smnorm{A}$ in the spin-1 model.

We take this limit separately at each length occurring in
Eq.~\eqref{sm:eq:closed}. Continuity preserves all three inequalities,
so they hold for $Z_L(\beta)$. We take this limit \emph{before}
iterating the doubling argument, which requires neither an error bound
uniform in $L$, nor eigenvector convergence, nor a single limiting
transfer operator. The finite slices need not form an exact Gibbs
semigroup; thermal squaring below uses only the exact exponential
after the limit.

\subsection{Torus identities and asymmetric doubling}
\label{sm:sec:asymmetric}

With the exact partition-function inequalities established, we now set
up the coupled doubling recurrence for the spatial and thermal purities:
\begin{equation}
 \Ps(n,\beta)=\frac{Z_{2n}(\beta)}{Z_n(\beta)^2},
 \qquad
 \Pt(L,\beta)=\frac{Z_L(2\beta)}{Z_L(\beta)^2}.
 \label{sm:eq:purities}
\end{equation}
All partition functions are strictly positive. The thermal purity is
the purity of the Gibbs state
$\rho_L(\beta)=e^{-\beta H_L}/Z_L(\beta)$.
The limiting spatial inequalities and the Boltzmann power sums give
$0<\Ps,\Pt\le1$ at even lengths and
\begin{align}
 1-\Ps(2n,\beta)&\le f(1-\Ps(n,\beta)), \notag\\
 1-\Pt(L,2\beta)&\le f(1-\Pt(L,\beta)).
 \label{sm:eq:two-squares}
\end{align}
For the spatial bound, divide the second limiting inequality in
Eq.~\eqref{sm:eq:closed} by $2Z_{2n}(\beta)^2$.
For the thermal bound, apply Eq.~\eqref{sm:eq:squaring} to the
normalized Boltzmann weights.

A torus of width $n$ and height $\beta$ represents $Z_n(\beta)$.
Gluing four copies in either order gives the exact cancellation
\begin{equation}
 \Ps(n,2\beta)\Pt(n,\beta)^2
 =\frac{Z_{2n}(2\beta)}{Z_n(\beta)^4}
 =\Ps(n,\beta)^2\Pt(2n,\beta).
 \label{sm:eq:torus}
\end{equation}
Solving for one factor and repeating at length $2n$ yields
\begin{align}
 \Ps(n,2\beta)
 &=\frac{\Ps(n,\beta)^2\Pt(2n,\beta)}{\Pt(n,\beta)^2}, \notag\\
 \Pt(4n,\beta)
 &=\frac{\Pt(2n,\beta)^2\Ps(2n,2\beta)}{\Ps(2n,\beta)^2}.
 \label{sm:eq:coupling}
\end{align}
Both identities follow directly by substituting
Eq.~\eqref{sm:eq:purities}; they concern exact partition functions.

Let $n\ge2$ be even, and let $0\le p,q<1$ bound the seed defects:
\begin{equation}
 \Ps(n,\beta)\ge1-p,\qquad
 \Pt(2n,\beta)\ge1-q.
 \label{sm:eq:seed-pair}
\end{equation}
To double both $n$ and $\beta$, we alternate the coupling identity with
squaring in space and imaginary time, feeding the updated spatial defect
into the thermal step. Write $z(a,b)=1-(1-a)^2(1-b)$.
The asymmetric update is
\begin{equation}
 p'=f(z(p,q)),\qquad q'=f(z(q,p')).
 \label{sm:eq:update}
\end{equation}
Assuming $p'<1$, its four steps are:
\begin{enumerate}
 \item The first identity in Eq.~\eqref{sm:eq:coupling}, with
 $\Pt(n,\beta)\le1$, gives
 $\Ps(n,2\beta)\ge(1-p)^2(1-q)=1-z(p,q)$.
 \item Spatial squaring and monotonicity of $f$ give
 $\Ps(2n,2\beta)\ge1-p'$.
 \item The second identity in Eq.~\eqref{sm:eq:coupling}, with
 $\Ps(2n,\beta)\le1$, gives
 $\Pt(4n,\beta)\ge(1-q)^2(1-p')=1-z(q,p')$.
 \item Thermal squaring gives $\Pt(4n,2\beta)\ge1-q'$.
\end{enumerate}
Thus $(p',q')$ bounds the defects at $(2n,2\beta)$.
Any upper bound below one substituted for $p'$ in the third step remains
valid; we use this rounding rule next.

\subsection{Six exact rational steps}
\label{sm:sec:orbit}

The analytic tail estimate in Sec.~\ref{sm:sec:gap} requires both
defects to be at most $1/100$. Because the initial spin-1 seed defects at
$(n_0,\beta_0)=(36,49/4)$ are larger, we first iterate the asymmetric
update for six steps in exact rational arithmetic.
The finite-size input is
\begin{equation}
 \Ps(36,49/4)\ge1-\frac{47916}{10^6},
 \qquad
 \Pt(72,49/4)\ge1-\frac{181521}{10^6}.
 \label{sm:eq:initial}
\end{equation}
Ref.~\cite{OpenAI2026Periodic} provides these certified seed bounds for
the periodic spin-1 chain; their finite-size spectral computation is not
repeated here. The following calculation verifies their propagation
under the recurrence.

All terminating decimals defining bounds denote exact rationals.
Use the ceiling grid $D^{-1}=10^{-7}$, $D=10^7$, and update
sequentially:
\begin{align}
 p_{j+1}&=\frac{\left\lceil Df(z(p_j,q_j))\right\rceil}{D},
 \notag\\
 q_{j+1}&=\frac{\left\lceil Df(z(q_j,p_{j+1}))\right\rceil}{D}.
 \label{sm:eq:rounding}
\end{align}
We round the spatial bound upward before the thermal update.
Table~\ref{sm:tab:orbit} lists the resulting orbit.

\begin{table}[ht]
 \centering
 \caption{Upward-rounded defect bounds at
 $n_j=36\,2^j$, $\beta_j=(49/4)2^j$.}
 \label{sm:tab:orbit}
 \begin{tabular}{rrrrr}
  \toprule
  $j$&$n_j$&$\beta_j$&$p_j$&$q_j$\\
  \midrule
  0&36&$49/4$&0.0479160&0.1815210\\
  1&72&$49/2$&0.0605003&0.1733838\\
  2&144&49&0.0686630&0.1632461\\
  3&288&98&0.0713693&0.1447318\\
  4&576&196&0.0633146&0.1055663\\
  5&1152&392&0.0376144&0.0446518\\
  6&2304&784&0.0084709&0.0055149\\
  \bottomrule
 \end{tabular}
\end{table}

To check the table in integer arithmetic, write $a_j=Dp_j$,
$b_j=Dq_j$, and set
\begin{equation}
 V_j=(D-a_j)^2(D-b_j),\qquad
 W_j=(D-b_j)^2(D-a_{j+1}).
 \label{sm:eq:integer-products}
\end{equation}
For each $j=0,\ldots,5$, the entries in
Table~\ref{sm:tab:orbit} satisfy
\begin{align}
 2(a_{j+1}-1)V_j^2
 &<D(D^3-V_j)^2\le2a_{j+1}V_j^2, \notag\\
 2(b_{j+1}-1)W_j^2
 &<D(D^3-W_j)^2\le2b_{j+1}W_j^2.
 \label{sm:eq:integer-check}
\end{align}
Multiplying the ceiling comparisons in Eq.~\eqref{sm:eq:rounding} by
their positive denominators gives Eq.~\eqref{sm:eq:integer-check}, since
$f(z(p_j,q_j))=(D^3-V_j)^2/(2V_j^2)$.
The right-hand inequalities alone suffice to verify the upper bounds.
Because all entries lie below one, the four-step update applies at every
row and yields
\begin{align}
 \Ps(2304,784)&\ge1-0.0084709>0.99,\notag \\
 \Pt(4608,784)&\ge1-0.0055149>0.99.
 \label{sm:eq:terminal}
\end{align}
Both defects at $(2304,784)$ now lie below the $1/100$ threshold.

\subsection{Scalar tail, interpolation, and the strict gap}
\label{sm:sec:gap}

Once both defects are at most $1/100$, a single scalar recurrence bounds
all subsequent doubling steps. For $0\le\varepsilon\le1/100$,
\begin{equation}
 f(1-(1-\varepsilon)^3)
 \le\frac{9\varepsilon^2}{2(1-3\varepsilon)^2}
 \le5\varepsilon^2\le\varepsilon.
 \label{sm:eq:scalar}
\end{equation}
The first inequality uses
$1-(1-\varepsilon)^3\le3\varepsilon$ and monotonicity of $f$;
the second uses $9/[2(97/100)^2]<5$.
If both input defects are at most $\varepsilon$, the spatial update
is therefore at most $5\varepsilon^2\le\varepsilon$.
Consequently, the thermal coupling input is also at most
$1-(1-\varepsilon)^3$, and its squared defect is at most
$5\varepsilon^2$.

Starting from Eq.~\eqref{sm:eq:terminal}, reset the scale index and put
\begin{equation}
 n_j=2304\,2^j,\qquad
 \beta_j=784\,2^j,\qquad
 \varepsilon_j=\frac{20^{-2^j}}5,\qquad j\ge0.
 \label{sm:eq:tail-scales}
\end{equation}
Here $\varepsilon_0=1/100$ and
$\varepsilon_{j+1}=5\varepsilon_j^2$.
Induction using Eq.~\eqref{sm:eq:scalar} proves
\begin{equation}
 \Ps(n_j,\beta_j)\ge1-\varepsilon_j,\qquad
 \Pt(2n_j,\beta_j)\ge1-\varepsilon_j.
 \label{sm:eq:tail}
\end{equation}

To cover every even length $L\ge2304$ between the dyadic scales $n_j$,
choose $j$ with $n_j\le L\le2n_j$. The limiting power-sum inequality in
Eq.~\eqref{sm:eq:closed}, applied at $\beta_j$ and $2\beta_j$, gives
\begin{align}
 Z_L(\beta_j)^2&\le Z_{n_j}(\beta_j)^{2L/n_j}, \notag\\
 Z_L(2\beta_j)&\ge Z_{2n_j}(2\beta_j)^{L/(2n_j)}.
 \label{sm:eq:interpolation-input}
\end{align}
Set $r=L/(2n_j)\in[1/2,1]$ and
\[
 x=\frac{Z_{2n_j}(2\beta_j)}{Z_{n_j}(\beta_j)^4}
   =\Pt(2n_j,\beta_j)\Ps(n_j,\beta_j)^2.
\]
Since $(1-\varepsilon_j)^3\le x\le1$, dividing the second inequality in
Eq.~\eqref{sm:eq:interpolation-input} by the first yields
\begin{equation}
 \Pt(L,\beta_j)\ge x^r\ge x
 \ge(1-\varepsilon_j)^3\ge1-3\varepsilon_j>\frac12.
 \label{sm:eq:interpolation}
\end{equation}

Finally, we convert this thermal purity bound at $\beta_j$ into a lower
bound on the spectral gap. Let $E_0(L)\le E_1(L)\le\cdots$ include
multiplicities, and let $w_m=e^{-\beta_j E_m(L)}/Z_L(\beta_j)$.
The largest Gibbs weight satisfies
$\Pt(L,\beta_j)=\sum_m w_m^2\le w_0$.
Equation~\eqref{sm:eq:interpolation} implies $w_0>1/2$, which
proves that the ground state is unique.
Moreover, $w_1\le1-w_0$. With
$\Delta_L=E_1(L)-E_0(L)$, we obtain
\begin{align}
 e^{-\beta_j\Delta_L}
 &=\frac{w_1}{w_0}
 \le\frac{3\varepsilon_j}{1-3\varepsilon_j} \notag\\
 &\le\frac{60}{97}\,20^{-2^j}
 <20^{-2^j},
 \label{sm:eq:strict-ratio}
\end{align}
where we used $1-3\varepsilon_j\ge97/100$.
Taking logarithms and dividing by $\beta_j=784\,2^j$ at this finite,
length-dependent temperature gives
\begin{equation}
 \boxed{
 \Delta_L>\frac{\log20}{784}
 \quad\text{for every even }L\ge2304
 }.
 \label{sm:eq:gap}
\end{equation}
The strict inequality comes from the finite prefactor $60/97<1$,
without taking a fixed-length limit $\beta\to\infty$.
Given the two seed bounds in Eq.~\eqref{sm:eq:initial}, the finite-matrix
transfer inequalities, rational orbit, scalar tail, and dyadic
interpolation establish Eq.~\eqref{sm:eq:gap}.
